% !TeX root = ../../../../../main.tex
% chapters/sections/chapter5/subsections/a_shock/summary.tex

\subsection{まとめ}
\label{sec:chapter5_a_shock_summary}
本節の締めくくりとして、供給ショック（\(a\)ショック）における各金融政策のパフォーマンスの違いをもたらした要因を総括し、あわせて名目総消費水準目標（NCLT）がここでも最上位クラスの健闘を見せた理由について考察する。

需要ショック（\(\beta\)ショック）時とは異なり、生産性の低下をもたらす供給ショックにおいては、「インフレ圧力」と「実体経済（消費・生産）の下押し圧力」が同時に発生する（スタグフレーション）。この環境下において、各政策は以下の厳しいトレードオフに対する姿勢の違いにより明暗が分かれることとなった。

\begin{enumerate}
    \item
    \label{itm:chapter5_a_shock_summary_real_economy_support}
        \textbf{インフレと価格分散を許容した実体経済の死守（上位群の勝因）}：
        総消費水準目標（CLT）や生産水準目標（OLT）、そして名目総消費水準目標（NCLT）や名目GDP水準目標（NGDPLT）などの上位群は、実質変数や名目支出の水準を維持しようとしたため、初期の強力なインフレとそれに伴う価格分散（\(\Delta_{t-1}^H\)）の大幅な増大を甘受した。しかし供給ショックにおいては、価格分散の蓄積による生産効率悪化のマイナス効果よりも、深刻な不況（消費と生産の大暴落）を防いだことのプラス効果の方がはるかに絶大であった。結果として、インフレを許容してでも実体経済を強力に下支えしたこれらの政策が、高い厚生評価を獲得することとなった。
    \item
    \label{itm:chapter5_a_shock_summary_inflation_control}
        \textbf{物価安定の過度な追求がもたらす実体経済の崩壊（下位群の敗因）}：
        対照的に、生産者物価インフレ目標（IT-PPI）や生産者物価水準目標（PPLT）などの下位群は、物価の安定を至上命題とした。これらの政策は、供給ショックによるインフレ圧力を機械的な金融引き締めによって力ずくでねじ伏せ、価格分散の発生を完全に抑え込むことには成功した。しかしその代償として、実体経済への強烈な下押し圧力が働き、消費と生産に致命的な大暴落をもたらした。消費者物価を目標とするIT-CPIやCPLT、各種テイラールール（TR-CPI, TR-PPI）なども同様に実体経済を過度に冷え込ませており、価格分散抑制のメリットを完全に相殺して余りあるほどの巨大な厚生損失を招く結果となった。
\end{enumerate}

\paragraph{名目総消費水準目標（NCLT）が健闘した理由とその頑健性}
\label{sec:chapter5_a_shock_summary_nclt_robustness}
上記のようなスタグフレーションの状況下において、実質変数のみを直接の目標とするCLTやOLTが最も高い厚生を達成したのは、インフレを完全に無視してでも実体経済を支え切ったためである。
一方で、名目変数を目標とするNCLTも、これら実質目標組に極めて肉薄する最上位クラスの厚生を記録し、見事な健闘を見せた。この理由は、NCLTが持つ「名目枠（物価水準 \(\times\) 実質消費）」という独自の構造にある。

供給ショックにより物価が上昇した際、物価のみを目標とする政策は直ちに強力な引き締めへと動く。しかしNCLTの場合、物価が上昇しても同時に実質消費が下落の危機にあれば、双方の変動が相殺し合い名目総消費の変動は小さくなる。そのため、中央銀行は実体経済を崩壊させるような過度な利上げ（引き締め）を行う必要がなくなり、結果的に実質目標組に近い「実体経済の強力な下支え」と「マイルドなインフレ許容」の最適なバランスを自動的に達成することができたのである。

前節の需要ショック（\(\beta\)ショック）においては、CLTがインフレの野放し（物価ドリフト）という弱点を露呈して評価を落としたのに対し、NCLTは名目アンカーとして機能し単独トップの最高評価を得ていた。
これと今回の供給ショック（\(a\)ショック）の結果を総合すると、NCLTは「需要ショックに対しては最強の回復力と物価安定をもたらし、中央銀行にとって対応が最も困難な供給ショックに対しても、実体経済の崩壊を防ぎ極めて傷の浅い対応ができる」という、あらゆるショックの源泉に対して極めて頑健（ロバスト）で優れた金融政策ルールであると結論付けることができる。